%% ma: L12-B14-0029
%% lop: 12
%% chuong: 5
%% bai: 14
%% dang: 12.14.17
%% loai: TN4
%% muc: TH
%% dap_an: "C"
%% nguon: "CĐ2.5.pdf giáo viên gửi 27/10/2026 (M07, Lê Minh Khôi – Nguyễn Anh Khôi, Chương II Chủ đề 5 Tích có hướng của hai vectơ), Câu 13"
%% kien_thuc: "Ba vectơ đồng phẳng, tích hỗn tạp (chương trình cũ, không có trong SGK KNTT 12)"
%% trang_thai: cach-ly
%% ly_do: "Vượt chương trình: điều kiện đồng phẳng của ba vectơ (tích hỗn tạp) không có trong SGK KNTT 12. Gợi ý: bỏ, hoặc đổi thành câu bốn điểm đồng phẳng giải bằng phương trình mặt phẳng."
%% ghi_chu: "Đề gốc đánh nhãn phương án cuối là C (trùng); đã đổi thành D"
\begin{ex}
Trong không gian $Oxyz$, cho ba vectơ $\vec a=(1;2;-1)$, $\vec b=(3;-1;0)$, $\vec c=(1;-5;2)$. Mệnh đề nào sau đây đúng?
\choice{$\vec a$ cùng phương với $\vec b$}{$\vec a$, $\vec b$, $\vec c$ không đồng phẳng}{\True $\vec a$, $\vec b$, $\vec c$ đồng phẳng}{$\vec a$ vuông góc với $\vec b$}
\loigiai{$\left[\vec a,\vec b\right]=(-1;-3;-7)$, $\left[\vec a,\vec b\right]\cdot\vec c=-1+15-14=0$ nên ba vectơ đồng phẳng; $\vec a\cdot\vec b=1\ne0$.}
\end{ex}
